\begingroup
\setlength{\emergencystretch}{2em}
\capituloTesis{Conclusiones}

\seccionTesis{Conclusión del objetivo general}

La evaluación de las veinte unidades estructura--sector del corredor EMP.LI105 Pedregal--Censo--Cumbray permitió caracterizar cuatro asociaciones entre la condición del drenaje vial y el desempeño funcional del pavimento, sin que los resultados sostuvieran una relación única para ambas variables consideradas en conjunto. Aunque las correlaciones convencionales de Spearman fueron de 0,5940 para CL--PCI, $-$0,6105 para CL--IRI, 0,1714 para CT--PCI y $-$0,1053 para CT--IRI, después del ajuste por tendencia espacial y de la corrección familiar de Holm únicamente CL--PCI conservó significancia estadística ($\rho$ ajustado = 0,5773; $p$ de Holm = 0,0180). En consecuencia, el objetivo general se atendió mediante las cuatro relaciones operacionalizadas, pero la evidencia fue específica y parcial, puesto que no se definió ni se calculó un contraste global de los dos constructos multidimensionales.

\seccionTesis{Conclusión del primer objetivo específico}

La continuidad longitudinal libre presentó una asociación positiva con la condición superficial del pavimento, expresada mediante un coeficiente convencional de Spearman de 0,5940 y una correlación residual ajustada de 0,5773; esta última alcanzó un valor $p$ espacial de 0,0045 y conservó significancia después de Holm ($p$ = 0,0180). Asimismo, al excluir sucesivamente cada unidad estructura--sector, los coeficientes permanecieron positivos entre 0,5316 y 0,6807, de manera que la dirección observada no dependió de una única unidad. Por consiguiente, el primer objetivo específico obtuvo respaldo estadístico para una asociación directa entre CL y PCI dentro del conjunto examinado, sin que ello permita atribuir causalidad a la continuidad del drenaje.

\seccionTesis{Conclusión del segundo objetivo específico}

La relación entre continuidad longitudinal libre y rugosidad mostró una dirección inversa, con $\rho$ de Spearman de $-$0,6105 y coeficiente residual de $-$0,5893, lo que describe una tendencia a encontrar menores valores de IRI donde la continuidad longitudinal fue mayor. No obstante, aunque el contraste espacial produjo $p$ = 0,0356, el valor corregido por Holm ascendió a 0,1068 y superó el umbral familiar de 0,05; por esta razón, el segundo objetivo específico identificó una correlación inversa de magnitud apreciable, pero la hipótesis H2 no alcanzó respaldo inferencial con la regla de decisión establecida. La falta de significancia corregida no constituye evidencia de inexistencia de esa relación.

\seccionTesis{Conclusión del tercer objetivo específico}

La abertura transversal libre registró una correlación positiva débil con el PCI ($\rho$ de Spearman = 0,1714), cuya magnitud aumentó a 0,2893 después del ajuste por coordenadas, sin alcanzar significancia en la aleatorización espacial ($p$ = 0,2675) ni tras la corrección de Holm ($p$ = 0,5350). Por ello, el tercer objetivo específico permitió estimar la dirección y la magnitud de CT--PCI en las veinte unidades analizadas, pero no aportó evidencia estadística suficiente para respaldar H3. Este resultado se refiere a la proporción geométrica de abertura libre observada en las estructuras transversales y no equivale a una medición de capacidad hidráulica de los pases de agua.

\seccionTesis{Conclusión del cuarto objetivo específico}

La asociación entre abertura transversal libre e índice internacional de rugosidad fue inversa y de baja magnitud, con $\rho$ de Spearman de $-$0,1053 y correlación residual ajustada de $-$0,1400, acompañada de valores $p$ espacial de 0,5217 y $p$ de Holm de 0,5350. El análisis por exclusión de una unidad produjo coeficientes entre $-$0,2211 y 0,0439, intervalo que comprendió ambos signos y mostró que la dirección del coeficiente no permaneció estable ante todas las exclusiones. Así, el cuarto objetivo específico quedó descrito cuantitativamente, aunque no se obtuvo respaldo inferencial para H4 ni evidencia de una relación consistente entre CT e IRI en el conjunto examinado.

\endgroup
